\documentclass[../DoAn.tex]{subfiles}
\begin{document}

\section{Giải pháp multi-tenant RAG chatbot cho tư vấn bán hàng nội thất}
\label{section:5.1}

Bài toán trung tâm của đồ án là xây dựng một hệ thống chatbot hỗ trợ tư vấn bán hàng nội thất cho nhiều cửa hàng trên cùng một nền tảng. Hệ thống không chỉ cần trả lời bằng ngôn ngữ tự nhiên, mà còn phải bảo đảm mỗi cửa hàng có dữ liệu, chatbot, hội thoại, lead và yêu cầu mua hàng riêng. Nếu dùng một chatbot chung cho tất cả cửa hàng, hệ thống dễ trả lời lẫn dữ liệu giữa các cửa hàng. Nếu triển khai riêng một hệ thống cho từng cửa hàng, chi phí vận hành và bảo trì sẽ tăng nhanh khi số lượng cửa hàng mở rộng.

Giải pháp được đề xuất là kết hợp mô hình multi-tenant với RAG chatbot. Trong đó, mỗi cửa hàng được biểu diễn như một tenant. Các thực thể quan trọng như chatbot, nguồn dữ liệu knowledge base, hội thoại, lead và yêu cầu mua hàng đều được gắn với tenant tương ứng. Backend chịu trách nhiệm xác thực, phân quyền, kiểm soát tenant và xử lý nghiệp vụ. AI Service chịu trách nhiệm truy xuất knowledge base, chuẩn bị ngữ cảnh và gọi Claude Sonnet để sinh câu trả lời tư vấn.

Khi người dùng gửi tin nhắn đến chatbot của một cửa hàng, hệ thống xác định chatbot và tenant tương ứng, sau đó chuyển yêu cầu sang AI Service. AI Service truy xuất dữ liệu liên quan từ knowledge base của tenant đó và sử dụng dữ liệu truy xuất làm ngữ cảnh để sinh câu trả lời. Cách xử lý này giúp chatbot trả lời dựa trên dữ liệu của đúng cửa hàng, thay vì chỉ dựa vào kiến thức chung của mô hình ngôn ngữ.

Một điểm quan trọng của giải pháp là liên kết hội thoại với nghiệp vụ bán hàng. Trong quá trình tư vấn, khi người dùng thể hiện ý định mua hàng hoặc cung cấp thông tin liên hệ, hệ thống có thể tạo lead hoặc yêu cầu mua hàng từ hội thoại. Nhân viên cửa hàng có thể xem lại nội dung hội thoại, nhận xử lý, phản hồi khách hàng và cập nhật trạng thái. Nhờ vậy, chatbot không chỉ là công cụ trả lời tự động mà còn là điểm đầu vào cho quy trình bán hàng của cửa hàng.

Kết quả đạt được là hệ thống có thể kết hợp ba thành phần chính: multi-tenant, RAG chatbot và xử lý lead/yêu cầu mua hàng. Hệ thống phục vụ được nhiều cửa hàng trên cùng một nền tảng, tư vấn dựa trên dữ liệu riêng của từng cửa hàng và chuyển một phần kết quả hội thoại thành thông tin nghiệp vụ để cửa hàng tiếp tục xử lý.

\section{Giải pháp phân tách dữ liệu tenant và dữ liệu tham chiếu công khai}
\label{section:5.2}

Trong hệ thống multi-tenant, dữ liệu bán hàng theo cửa hàng phải được cô lập nghiêm ngặt theo tenant. Chatbot của một cửa hàng không được sử dụng sản phẩm, chính sách hoặc hội thoại của cửa hàng khác để tư vấn. Tuy nhiên, hệ thống cũng có các chức năng không đại diện cho một cửa hàng cụ thể, chẳng hạn tư vấn tổng quát, so sánh chung và tham chiếu giá sản phẩm. Các chức năng này cần nguồn dữ liệu rộng hơn dữ liệu của một tenant, nhưng không được làm mất nguyên tắc tách biệt dữ liệu.

Giải pháp được đề xuất là phân tách dữ liệu thành hai phạm vi: dữ liệu theo tenant và dữ liệu tham chiếu công khai. Dữ liệu theo tenant là dữ liệu phục vụ trực tiếp cho một cửa hàng, bao gồm sản phẩm, chính sách, bài viết và nguồn tri thức do cửa hàng cấu hình. Dữ liệu này luôn gắn với \texttt{tenant\_id} và chỉ được sử dụng trong các luồng thuộc tenant tương ứng.

Dữ liệu tham chiếu công khai là dữ liệu được thu thập từ các nguồn công khai, dùng cho các chế độ không đại diện cho một cửa hàng cụ thể như \texttt{general\_compare} và \texttt{market\_price}. Dữ liệu này được xem là kho tham chiếu độc lập, không phải dữ liệu riêng của nhiều tenant bị trộn lại. Mỗi bản ghi tham chiếu cần giữ thông tin nguồn như tên cửa hàng, URL và thời điểm thu thập để người dùng hiểu phạm vi tham khảo của kết quả.

Trong pipeline dữ liệu, mỗi nguồn dữ liệu được gắn với trường \texttt{visibility}. Các giá trị chính gồm \texttt{tenant\_private}, \texttt{tenant\_public} và \texttt{public\_reference}. Dữ liệu \texttt{tenant\_private} và \texttt{tenant\_public} phải có \texttt{tenant\_id}; dữ liệu \texttt{public\_reference} không được dùng như dữ liệu riêng của tenant. Ở mức thư mục, dữ liệu tenant được đưa vào \texttt{output/tenants/\{tenant\_id\}}, còn dữ liệu tham chiếu được đưa vào \texttt{output/reference}.

Kết quả của giải pháp là hệ thống có cách tổ chức dữ liệu rõ ràng hơn. Chat bán hàng theo cửa hàng chỉ sử dụng dữ liệu của tenant tương ứng, trong khi chat tổng quát và tham chiếu giá sử dụng kho tham chiếu công khai. Cách thiết kế này giúp bảo đảm tenant isolation, đồng thời vẫn hỗ trợ được các chức năng so sánh và tham chiếu giá trong phạm vi dữ liệu hệ thống.

\section{Giải pháp kiểm soát ba chế độ hội thoại}
\label{section:5.3}

Trong hệ thống chatbot nội thất, không phải mọi hội thoại đều có cùng mục tiêu. Một người dùng có thể chat với chatbot của một cửa hàng để hỏi mua sản phẩm, nhưng cũng có thể hỏi tư vấn chung hoặc tham khảo mức giá. Nếu các luồng này được xử lý giống nhau, hệ thống dễ phát sinh hành vi sai, chẳng hạn tạo yêu cầu mua hàng trong một cuộc hội thoại chỉ dùng để tham chiếu giá.

Giải pháp được đề xuất là chuẩn hóa ba chế độ hội thoại: \texttt{tenant\_sales}, \texttt{general\_compare} và \texttt{market\_price}. Chế độ \texttt{tenant\_sales} phục vụ tư vấn bán hàng theo một cửa hàng cụ thể. Chế độ này sử dụng dữ liệu của tenant tương ứng, có thể hỏi làm rõ nhu cầu, tư vấn sản phẩm và tạo lead hoặc yêu cầu mua hàng khi đủ điều kiện.

Chế độ \texttt{general\_compare} phục vụ tư vấn và so sánh chung. Chế độ này sử dụng dữ liệu tham chiếu công khai, giúp người dùng hiểu các lựa chọn sản phẩm và tiêu chí so sánh ở mức tham khảo. Chế độ này không tạo lead hoặc yêu cầu mua hàng. Chế độ \texttt{market\_price} phục vụ tham chiếu giá sản phẩm, sử dụng dữ liệu tham chiếu công khai để đưa ra khoảng giá hoặc nhận xét mức giá trong phạm vi dữ liệu hệ thống. Chế độ này cũng không tạo yêu cầu mua hàng và không thay thế báo giá thị trường đầy đủ.

Mode hội thoại được gửi kèm trong contract giữa backend và AI Service. AI Service dùng mode để chọn prompt, phạm vi dữ liệu và logic phản hồi phù hợp. Backend vẫn là lớp kiểm soát cuối cùng đối với các tác động nghiệp vụ. Backend chỉ tạo yêu cầu mua hàng khi mode yêu cầu là \texttt{tenant\_sales}, mode cuối cùng cũng là \texttt{tenant\_sales}, hội thoại thuộc tenant hợp lệ và các điều kiện nghiệp vụ khác được thỏa mãn.

Kết quả của giải pháp là hệ thống phân biệt rõ ba nhóm nhu cầu hội thoại: tư vấn bán hàng theo cửa hàng, tư vấn/so sánh chung và tham chiếu giá. Việc chuẩn hóa mode giúp giảm tình trạng lẫn hành vi giữa các luồng, đặc biệt bảo đảm \texttt{general\_compare} và \texttt{market\_price} không tạo yêu cầu mua hàng ngoài ý muốn.

\section{Giải pháp pipeline dữ liệu cho catalog sản phẩm và knowledge base}
\label{section:5.4}

Chất lượng câu trả lời của chatbot phụ thuộc trực tiếp vào chất lượng dữ liệu đầu vào. Với sản phẩm nội thất, dữ liệu cần có các thông tin như tên sản phẩm, giá, chất liệu, kích thước, loại sản phẩm, phong cách, mô tả, cửa hàng và URL nguồn. Ngoài dữ liệu sản phẩm, chatbot theo cửa hàng còn cần các tài liệu như chính sách bảo hành, vận chuyển, giới thiệu cửa hàng hoặc bài viết hướng dẫn chọn sản phẩm.

Giải pháp được đề xuất là xây dựng thư mục \texttt{data\_pipeline} độc lập với runtime chatbot. Thư mục này chứa file CSV đầu vào, script crawl, parser, script normalize và script validate. Nguồn dữ liệu được khai báo bằng CSV để người phát triển kiểm soát rõ crawler sẽ truy cập URL nào, tránh việc quét tự động toàn bộ website một cách khó kiểm soát.

Dữ liệu sau crawl được chia thành hai nhóm đầu ra. Nhóm thứ nhất là \texttt{products.clean.jsonl}, chứa các bản ghi sản phẩm có cấu trúc như tên, loại, giá, chất liệu, kích thước, phong cách, mô tả, ảnh, URL nguồn và thời điểm crawl. Nhóm thứ hai là \texttt{knowledge\_docs.clean.jsonl}, chứa các tài liệu dạng chính sách, giới thiệu hoặc bài viết, gồm tiêu đề, nội dung, loại tài liệu, URL nguồn và thời điểm crawl.

Pipeline có bước validate để kiểm tra chất lượng dữ liệu trước khi đưa vào chatbot. Với sản phẩm, hệ thống kiểm tra số lượng bản ghi, tỷ lệ có giá, tỷ lệ có chất liệu/kích thước, URL trùng và các trường thiếu. Với tài liệu knowledge, hệ thống kiểm tra tiêu đề, nội dung, nguồn và phạm vi tenant. Các lỗi về visibility và tenant cũng được kiểm tra để tránh đưa dữ liệu sai phạm vi vào runtime chatbot.

Nguyên tắc quan trọng của pipeline là không tự tạo dữ liệu không có căn cứ. Nếu website không cung cấp giá, chất liệu hoặc kích thước, trường tương ứng được để \texttt{null}. Pipeline giữ lại raw text và source URL để truy vết. Kết quả đạt được là hệ thống có quy trình chuẩn bị dữ liệu rõ ràng hơn so với nhập tay rời rạc hoặc dùng dữ liệu mock, đồng thời hỗ trợ tốt hơn cho các chức năng so sánh sản phẩm và tham chiếu giá.

\section{Giải pháp triển khai nhẹ với Claude API và local model fallback}
\label{section:5.5}

Một khó khăn khi triển khai chatbot dùng mô hình ngôn ngữ là yêu cầu tài nguyên. Nếu chạy local model trên VPS CPU-only, thời gian phản hồi có thể dài và bộ nhớ sử dụng cao. Nếu dùng GPU, chi phí triển khai tăng lên đáng kể. Trong phạm vi đồ án, mục tiêu là triển khai được hệ thống demo ổn định trên VPS phổ thông, phục vụ các luồng chính như tư vấn theo cửa hàng, so sánh chung và tham chiếu giá.

Giải pháp được lựa chọn là sử dụng Claude API làm provider chính trong môi trường triển khai, còn local model được giữ như phương án thử nghiệm hoặc fallback kỹ thuật. API key của Claude được cấu hình bằng biến môi trường, không lưu trong mã nguồn, tài liệu hoặc cơ sở dữ liệu. Cách làm này giúp bảo vệ khóa tích hợp và dễ thay đổi key khi cần.

Trong Docker Compose, hệ thống gồm các service chính: \texttt{app}, \texttt{chatbot-api} và \texttt{postgres}. Backend Spring Boot chạy trong service \texttt{app}; AI Service FastAPI chạy trong service \texttt{chatbot-api}; PostgreSQL chạy trong service \texttt{postgres}. Backend gọi AI Service qua địa chỉ nội bộ \texttt{http://chatbot-api:8000}. Cách tổ chức này giúp hệ thống không phụ thuộc vào đường dẫn Python cục bộ và dễ chạy lại trên VPS.

Để tránh local model tiêu tốn tài nguyên ngoài ý muốn, cấu hình mặc định không load local model khi khởi động. Local pipeline chỉ được bật khi có cấu hình rõ ràng. Với VPS, các thông tin triển khai được lưu trong file \texttt{.env}, file này không được commit lên Git. Các port nội bộ như PostgreSQL và AI Service không cần mở public; người dùng truy cập hệ thống qua backend hoặc reverse proxy.

Kết quả của giải pháp là hệ thống chạy nhẹ hơn trong môi trường demo. Các service chính được đóng gói bằng Docker Compose, backend gọi AI Service qua HTTP nội bộ, provider chính là Claude API và local model không được load mặc định. Cách triển khai này giảm rủi ro thiếu tài nguyên, giảm phụ thuộc vào môi trường máy phát triển và phù hợp hơn với mục tiêu demo đồ án.

\section{Tổng kết đóng góp}
\label{section:5.6}

Các giải pháp trong chương này tập trung vào các vấn đề cốt lõi của đồ án. Thứ nhất, hệ thống kết hợp multi-tenant với RAG chatbot để tư vấn bán hàng nội thất cho nhiều cửa hàng trên cùng một nền tảng. Thứ hai, hệ thống phân tách dữ liệu tenant và dữ liệu tham chiếu công khai để vừa bảo đảm tenant isolation, vừa hỗ trợ chat tổng quát và tham chiếu giá. Thứ ba, hệ thống chuẩn hóa ba chế độ hội thoại để kiểm soát phạm vi dữ liệu và tác động nghiệp vụ. Thứ tư, hệ thống xây dựng pipeline dữ liệu riêng để chuẩn bị catalog sản phẩm và knowledge base có kiểm soát. Cuối cùng, hệ thống được triển khai theo hướng nhẹ bằng Docker Compose và Claude API, giữ local model như phương án thử nghiệm thay vì phụ thuộc vào GPU.

Những đóng góp này tập trung vào một bài toán cụ thể: hỗ trợ tư vấn bán hàng nội thất trong môi trường nhiều cửa hàng, có dữ liệu riêng, có hội thoại nhiều lượt và có khả năng chuyển kết quả tư vấn thành thông tin nghiệp vụ để cửa hàng tiếp tục xử lý. Đây là điểm khác biệt chính của đồ án so với một chatbot hỏi đáp thông thường.

\end{document}